\section{РЕАЛИЗАЦИЯ И РЕЗУЛЬТАТЫ}

\subsection{Пользовательский интерфейс}

Графический интерфейс программного комплекса построен на базе Pygame и состоит из пяти основных экранов, образующих последовательный рабочий поток: главное меню, редактор карт, выбор алгоритма, визуализация симуляции и экраны батч-режима. Навигация между экранами осуществляется клавишей ESC (возврат на предыдущий экран) и кнопками действий.

\subsubsection{Главное меню}

Главное меню предоставляет выбор между тремя режимами работы:

\begin{itemize}
    \item \textbf{Create} --- создание нового сценария в редакторе карт;
    \item \textbf{Load} --- загрузка ранее сохранённого сценария из JSON-файла (через системный диалог);
    \item \textbf{Batch Mode} --- переход к автоматизированному пакетному тестированию.
\end{itemize}

% === МЕСТО ДЛЯ РИСУНКА ===
\begin{figure}
  \centering
  %\fbox{\parbox{0.6\textwidth}{\centering\vspace{3cm}\textit{Скриншот главного меню}\vspace{3cm}}}
  \includegraphics[width=0.6\textwidth]{inc/main_menu}

  \caption{Главное меню программного комплекса}
  \label{fig:main_menu}
\end{figure}

\subsubsection{Редактор карт}

Редактор карт позволяет интерактивно создавать и модифицировать сценарии. Левая часть экрана содержит сеточную карту, правая --- панель инструментов. Пользователь может:

\begin{itemize}
    \item рисовать и удалять препятствия щелчками мыши;
    \item размещать стартовые позиции агентов и целевые позиции;
    \item изменять размер сетки непосредственно в панели настроек;
    \item сохранять и загружать сценарии в формате JSON.
\end{itemize}

% === МЕСТО ДЛЯ РИСУНКА ===
\begin{figure}
  \centering
  %\fbox{\parbox{0.75\textwidth}{\centering\vspace{4cm}\textit{Скриншот редактора карт}\vspace{4cm}}}
  \includegraphics[width=0.6\textwidth]{inc/editor}
  \caption{Редактор карт с размещёнными агентами и целями}
  \label{fig:editor}
\end{figure}

\subsubsection{Экран выбора алгоритма}

На экране выбора алгоритма отображается предварительный просмотр карты и панель с двумя вкладками: MAPF и MRTA. Пользователь выбирает алгоритм из доступного списка, после чего отображаются его характеристики (тип архитектуры, режим работы). При выборе MRTA-алгоритма дополнительно предлагается выбрать MAPF-алгоритм в качестве навигационного планировщика. Перед запуском выполняется валидация: проверяется, что число агентов и целей совместимо с выбранным типом алгоритма.

\subsubsection{Визуализация симуляции}

Экран симуляции является центральным элементом интерактивного режима. Левая часть отображает сеточную карту с агентами, целями (или задачами в режиме MRTA), препятствиями и опционально --- построенными путями агентов. Правая боковая панель содержит:

\begin{itemize}
    \item информацию об алгоритме (имя, тип, режим);
    \item текущий статус (PLANNING, RUNNING, DONE, FAILED);
    \item метрики в реальном времени: Makespan, SOC, время планирования (общее, MRTA, MAPF), память;
    \item управление: пауза (Space), перезапуск (R), переключение отображения путей (P), регулировка скорости (+/--).
\end{itemize}

Агенты отображаются цветными квадратами с номерами, цели в режиме MAPF --- перекрестиями соответствующего цвета, задачи в режиме MRTA --- ромбами с номерами. Пути отображаются в виде цветных точек, позволяя визуально оценить маршруты и выявить потенциальные зоны скопления.

% === МЕСТО ДЛЯ РИСУНКА ===
\begin{figure}
  \centering
  %\fbox{\parbox{0.75\textwidth}{\centering\vspace{4.5cm}\textit{Скриншот симуляции с отображением путей и метрик}\vspace{4.5cm}}}
  \includegraphics[width=0.6\textwidth]{inc/MAPF}
  \caption{Визуализация симуляции (MAPF-режим, алгоритм CBS)}
  \label{fig:simulation}
\end{figure}

\subsubsection{Батч-режим}

Батч-режим реализован в виде последовательности из трёх экранов:

\begin{enumerate}[arabic]
    \item \textbf{Выбор алгоритмов} --- чекбоксы для MAPF- и MRTA-алгоритмов с отображением общего числа комбинаций;
    \item \textbf{Настройка параметров} --- размер сетки, диапазон числа роботов (min/max/step), плотность препятствий, число сценариев, стратегия размещения, таймаут, импорт карт;
    \item \textbf{Прогресс} --- индикатор выполнения в реальном времени (эксперименты выполняются в отдельном потоке), текущая комбинация, кнопка открытия папки с результатами.
\end{enumerate}

% === МЕСТО ДЛЯ РИСУНКА ===
\begin{figure}
  \centering
  %\fbox{\parbox{0.75\textwidth}{\centering\vspace{3.5cm}\textit{Скриншот экрана батч-режима (настройка параметров)}\vspace{3.5cm}}}
  \includegraphics[width=0.6\textwidth]{inc/batch}
  \caption{Экран настройки параметров пакетного тестирования}
  \label{fig:batch_settings}
\end{figure}

\subsection{Сравнительный анализ алгоритмов}

\subsubsection{Методика эксперимента}

Для объективного сравнения реализованных алгоритмов были проведены серии экспериментов с использованием батч-режима. Параметры экспериментов:

\begin{itemize}
    \item размер сетки: $12 \times 12$ клеток;
    \item плотность препятствий: 15\%;
    \item диапазон числа роботов: от 2 до 8 с шагом 2;
    \item стратегия размещения: uniform;
    \item число сценариев на каждое значение числа роботов: 3;
    \item таймаут на один прогон: 60 секунд;
    \item проверка достижимости: включена.
\end{itemize}

Для каждой комбинации (MRTA-алгоритм, MAPF-алгоритм, число роботов) измерялись: makespan, SOC, суммарное время планирования, пиковое потребление памяти и флаг успешного завершения. Результаты усреднялись по всем сценариям (seeds) для каждого значения числа роботов.

\subsubsection{Анализ MAPF-алгоритмов}

Результаты экспериментов позволяют выявить характерные особенности каждого класса MAPF-алгоритмов.

\textbf{Оптимальные алгоритмы (CBS, M*).} CBS обеспечивает гарантированно оптимальные решения по метрике SOC, однако время планирования растёт экспоненциально с увеличением числа агентов и плотности конфликтов. Вариант CBS с SIPP в качестве нижнеуровневого планировщика показывает ускорение при большом количестве ограничений за счёт более компактного представления пространства поиска. M* эффективен в сценариях с разреженными конфликтами, где большинство агентов не взаимодействуют: в таких случаях поиск остаётся фактически низкоразмерным. Однако при плотных конфликтах M* деградирует до полного совместного поиска.

\textbf{Субоптимальные алгоритмы (ECBS, EECBS).} ECBS и EECBS обеспечивают значительное ускорение по сравнению с CBS за счёт ослабления требования оптимальности. EECBS, благодаря более точной эвристике на верхнем уровне, как правило, находит решения ближе к оптимуму при том же бюджете времени. Эти алгоритмы представляют хороший компромисс между качеством решения и скоростью планирования для сценариев со средним числом агентов.

\textbf{Децентрализованные алгоритмы (CA*, WHCA*, PIBT).} CA* и WHCA* показывают наименьшее время планирования за счёт последовательной архитектуры, однако качество решений существенно зависит от порядка планирования агентов. WHCA* с ограниченным окном лучше справляется с динамическими сценариями, поскольку периодически перепланирует маршруты. PIBT, планирующий на один шаг вперёд, демонстрирует лучшую масштабируемость при большом числе агентов, но может генерировать неоптимальные пути с высоким SOC.

В таблице~\ref{tab:mapf_comparison} приведена качественная сравнительная характеристика MAPF-алгоритмов.

\begin{table}
    \caption{Сравнительная характеристика MAPF-алгоритмов}
    \begin{tabularx}{\textwidth}{|l|c|c|X|}\hline
    Алгоритм & Качество & Скорость & Область применения \\ \hline
    CBS (A*)  & отличное & низкая   & Задачи с малым числом агентов, где критична оптимальность \\ \hline
    CBS (SIPP)& отличное & средняя  & Аналогично CBS, эффективнее при плотных ограничениях \\ \hline
    ECBS      & хорошее  & средняя  & Среднее число агентов, допустима ограниченная субоптимальность \\ \hline
    EECBS     & хорошее  & выше ср. & Аналогично ECBS, с улучшенной эвристикой \\ \hline
    M*        & отличное & переменная & Разреженные конфликты, группы агентов с малым пересечением \\ \hline
    CA*       & среднее  & высокая  & Большое число агентов, критично время планирования \\ \hline
    WHCA*     & среднее  & высокая  & Динамические сценарии, большое число агентов \\ \hline
    PIBT      & ниже ср. & очень высокая & Реактивное управление, очень большие группы агентов \\ \hline
    \end{tabularx}
    \label{tab:mapf_comparison}
\end{table}

\subsubsection{Анализ MRTA-алгоритмов}

Анализ MRTA-алгоритмов выявляет различия между online- и queue-подходами.

\textbf{Online-алгоритмы (венгерский, жадный).} Венгерский алгоритм обеспечивает оптимальное решение задачи назначения на каждой итерации, что приводит к лучшему суммарному расстоянию при небольшом числе задач. Жадный алгоритм значительно быстрее, но может давать существенно худшие назначения, особенно когда расположение задач неоднородно. Недостатком online-подхода является отсутствие глобальной оптимизации: каждое назначение учитывает только текущее состояние.

\textbf{Queue-алгоритмы (последовательный аукцион, минимальная стоимость потока, комбинаторный аукцион).} Алгоритмы с построением цепочек учитывают все задачи одновременно, что позволяет минимизировать суммарные перемещения. Минимальная стоимость потока обеспечивает оптимальное решение на каждом слое назначений. Последовательный аукцион даёт хороший баланс скорости и качества. Комбинаторный аукцион с локальной оптимизацией обменом показывает лучшие результаты в сценариях, где существуют естественные кластеры задач, поскольку может назначить близкие задачи одному роботу.

Выбор алгоритма MRTA существенно влияет на итоговый SOC и makespan, однако зависит также от выбранного MAPF-планировщика: оптимальное назначение не гарантирует оптимальных маршрутов, и наоборот.

\subsubsection{Результаты пакетного тестирования}

На рисунках~\ref{fig:makespan_results}--\ref{fig:success_results} представлены графики зависимости ключевых метрик от числа роботов, полученные в ходе пакетного тестирования.

% === МЕСТО ДЛЯ ГРАФИКОВ ===
% Графики генерируются автоматически батч-режимом в папку batch_results/
% Скопируйте нужные файлы в report/inc/ и раскомментируйте \includegraphics

\begin{figure}
  \centering
  \fbox{\parbox{0.85\textwidth}{\centering\vspace{4cm}\textit{График Makespan vs Number of Robots}\vspace{4cm}}}
  % \includegraphics[width=0.85\textwidth]{inc/makespan_vs_robots}
  \caption{Зависимость makespan от числа роботов}
  \label{fig:makespan_results}
\end{figure}

\begin{figure}
  \centering
  \fbox{\parbox{0.85\textwidth}{\centering\vspace{4cm}\textit{График SOC vs Number of Robots}\vspace{4cm}}}
  % \includegraphics[width=0.85\textwidth]{inc/soc_vs_robots}
  \caption{Зависимость SOC от числа роботов}
  \label{fig:soc_results}
\end{figure}

\begin{figure}
  \centering
  \fbox{\parbox{0.85\textwidth}{\centering\vspace{4cm}\textit{График Planning Time vs Number of Robots}\vspace{4cm}}}
  % \includegraphics[width=0.85\textwidth]{inc/planning_time_vs_robots}
  \caption{Зависимость времени планирования от числа роботов}
  \label{fig:planning_results}
\end{figure}

\begin{figure}
  \centering
  \fbox{\parbox{0.85\textwidth}{\centering\vspace{4cm}\textit{График Success Rate vs Number of Robots}\vspace{4cm}}}
  % \includegraphics[width=0.85\textwidth]{inc/success_rate_vs_robots}
  \caption{Зависимость доли успешных прогонов от числа роботов}
  \label{fig:success_results}
\end{figure}

С увеличением числа роботов наблюдаются следующие закономерности:

\begin{enumerate}[arabic]
    \item \textbf{Makespan} растёт для всех алгоритмов, однако темп роста существенно различается: оптимальные алгоритмы (CBS) обеспечивают минимальный makespan, но при большом числе агентов могут не уложиться в таймаут; субоптимальные (ECBS, EECBS) дают близкие к оптимальному значения при значительно меньшем времени;
    \item \textbf{SOC} демонстрирует аналогичную тенденцию: децентрализованные алгоритмы (CA*, WHCA*) склонны к генерации избыточных путей, особенно при плотном расположении агентов;
    \item \textbf{Время планирования} является наиболее дифференцирующей метрикой: CBS и M* показывают экспоненциальный рост, ECBS/EECBS --- умеренный, CA*/WHCA*/PIBT --- близкий к линейному;
    \item \textbf{Success Rate} снижается с ростом числа агентов для всех алгоритмов из-за ограничений таймаута и сложности карт, но децентрализованные алгоритмы сохраняют высокую долю успешных прогонов.
\end{enumerate}

\subsection{Обсуждение результатов}

Проведённый анализ подтверждает, что не существует универсально лучшего алгоритма: выбор зависит от условий среды.

Для задач с малым числом агентов (2--4), где критична оптимальность, рекомендуется использование CBS или M*. Для сценариев со средним числом агентов (4--8) оптимальный баланс скорости и качества обеспечивают ECBS и EECBS. Для задач с большим числом агентов или жёсткими ограничениями на время планирования следует использовать PIBT или WHCA*.

При комбинировании MRTA и MAPF алгоритмы с построением цепочек (queue-режим) показывают преимущество при большом числе задач, поскольку минимизируют холостые перемещения роботов. Однако online-подход (венгерский алгоритм) более устойчив к непредвиденным ситуациям, поскольку перепланирует назначения по мере завершения задач.

Разработанный программный комплекс позволяет проводить подобный анализ для произвольных конфигураций среды, что даёт возможность обоснованно выбирать алгоритмы для конкретных прикладных задач.
